%! TEX program = xelatex
%-------------------------
% Chinese Resume in Latex
% Author : Chen Xiang
% Based off of: https://github.com/jakegut/resume
% License : MIT
%------------------------

\documentclass[letterpaper,11pt,fontset=none]{article}

\usepackage{fontspec}
\usepackage{xeCJK}
% \setCJKmainfont{Heiti SC}[
%   UprightFont = * Light,
%   BoldFont = * Medium
% ]
\setmainfont{SF Pro Text}[
  BoldFont = * Bold,
  FontFace = {l}{n}{*-Regular},
  FontFace = {sb}{n}{*-Medium}
]
\setsansfont{SF Pro Text}[
  BoldFont = * Bold,
  FontFace = {l}{n}{*-Regular},
  FontFace = {sb}{n}{*-Medium}
]
\setCJKmainfont{PingFang SC}[
  BoldFont = * Semibold,
  UprightFont = * Regular,
  FontFace = {l}{n}{*-Regular},
  FontFace = {sb}{n}{*-Medium}
]

\usepackage{latexsym}
\usepackage[empty]{fullpage}
\usepackage{titlesec}
\usepackage{marvosym}
\usepackage[usenames,dvipsnames]{color}
\usepackage{verbatim}
\usepackage{enumitem}
\usepackage[hidelinks]{hyperref}
\usepackage{fancyhdr} 
\usepackage[english]{babel}
\usepackage{tabularx}
\usepackage{xstring}
\usepackage{ifthen}
\usepackage{fontawesome5}
\usepackage[normalem]{ulem} % For customizable underlines
% \input{glyphtounicode}

%----------FONT OPTIONS----------
% sans-serif
% \usepackage[sfdefault]{FiraSans}
% \usepackage[sfdefault]{roboto}
% \usepackage[sfdefault]{noto-sans}
% \usepackage[default]{sourcesanspro}
% \usepackage[sfdefault]{carlito}

% serif
% \usepackage{CormorantGaramond}
% \usepackage{charter}


\pagestyle{fancy}
\fancyhf{} % clear all header and footer fields
\fancyfoot{}
\renewcommand{\headrulewidth}{0pt}
\renewcommand{\footrulewidth}{0pt}
\linespread{1.1} 

% Adjust margins
\addtolength{\oddsidemargin}{-0.55in}
\addtolength{\evensidemargin}{-0.55in}
\addtolength{\textwidth}{1in}
\addtolength{\topmargin}{-.5in}
\addtolength{\textheight}{1.2in}
% \addtolength{\footskip}{-.9in}

\urlstyle{same}
 
\raggedbottom
\raggedright
\setlength{\tabcolsep}{0in}

% Sections formatting
\titleformat{\section}{
  \vspace{-9pt}\raggedright
}{}{0em}{}[\color{black}\titlerule \vspace{-7pt}]

% Ensure that generate pdf is machine readable/ATS parsable
% \pdfgentounicode=1

% Phone number display control (default: true)
\providecommand{\showphone}{true}
\providecommand{\phonenumber}{+852 0000-0000}
\providecommand{\phonenumbertwo}{+86 000-0000-0000}

\newcommand{\sbseries}{\fontseries{sb}\selectfont}
\newcommand{\textsb}[1]{{\sbseries #1}}

%-------------------------
% Custom commands

% Custom thin underline command
\newcommand{\thinunderline}[1]{\uline{#1}}
% Adjust underline thickness (0.4pt is thinner than default)
\renewcommand{\ULthickness}{0.2pt}

\newcommand{\resumeItem}[1]{
  \item\fontseries{l}\fontsize{10pt}{11.5pt}\selectfont\color[rgb]{0.1,0.1,0.1}{
    {#1 \vspace{-2pt}}
  }
}

% \newcommand{\resumeSubheading}[4]{
%   \vspace{-2pt}\item
%     \begin{tabular*}{0.97\textwidth}[t]{l@{\extracolsep{\fill}}r}
%       \textbf{#1} & #2 \\
%       \textit{\small#3} & \textit{\small #4} \\
%     \end{tabular*}\vspace{-5pt}
% }

\newcommand{\resumeSubheading}[2]{
  \vspace{-2pt}\item
    \begin{tabular*}{0.97\textwidth}[t]{l@{\extracolsep{\fill}}r}
      \textbf{#1} & #2 \\
    \end{tabular*}\vspace{-9pt}
}


\newcommand{\resumeSubSubheading}[2]{
    \item
    \begin{tabular*}{0.97\textwidth}{l@{\extracolsep{\fill}}r}
      {\small\textsb{#1}} & {\footnotesize #2} \\
    \end{tabular*}\vspace{-9pt}
}

\newcommand{\resumeProjectHeading}[2]{
    \item
    \begin{tabular*}{0.97\textwidth}{l@{\extracolsep{\fill}}r}
      \small#1 & {\footnotesize #2} \\
    \end{tabular*}\vspace{-9pt}
}

\newcommand{\resumeSubItem}[1]{\resumeItem{#1}\vspace{-2pt}}

\renewcommand\labelitemii{$\vcenter{\hbox{\tiny$\bullet$}}$}

\newcommand{\resumeSubHeadingListStart}{\begin{itemize}[leftmargin=0.1in, label={}]}
\newcommand{\resumeSubHeadingListEnd}{\end{itemize}}
\newcommand{\resumeItemListStart}{\begin{itemize}[leftmargin=0.18in, rightmargin=0.015\textwidth]}
\newcommand{\resumeItemListEnd}{\end{itemize}\vspace{-9pt}}

%-------------------------------------------
%%%%%%  RESUME STARTS HERE  %%%%%%%%%%%%%%%%%%%%%%%%%%%%


\begin{document}

% MARK: HEADING
%----------HEADING----------
% \begin{tabular*}{\textwidth}{l@{\extracolsep{\fill}}r}
%   \textbf{\href{http://sourabhbajaj.com/}{\Large Sourabh Bajaj}} & Email : \href{mailto:sourabh@sourabhbajaj.com}{sourabh@sourabhbajaj.com}\\
%   \href{http://sourabhbajaj.com/}{http://www.sourabhbajaj.com} & Mobile : +1-123-456-7890 \\
% \end{tabular*}

\begin{center}
	\textbf{\huge 陈 想}\\ \vspace{2pt}
	\footnotesize
	\ifthenelse{\equal{\showphone}{true}}{{\raisebox{-0.5pt}\faPhone} \space \phonenumbertwo \space $|$}{}
	\raisebox{-1pt}{\faEnvelope} \space \href{mailto:xiiang.ch@gmail.com}{\thinunderline{xiiang.ch@gmail.com}} $|$
	\raisebox{-1pt}{\faGlobe} \space \href{https://cxiang.site}{\thinunderline{cxiang.site}} $|$
	\raisebox{-1pt}{\faGithub} \space \href{https://github.com/Xiang-CH}{\thinunderline{github.com/Xiang-CH}} $|$
  \raisebox{-0.5pt}{\faMapMarker*} \space \scriptsize 上海/香港
\end{center}

% MARK: EDUCATION
%-----------EDUCATION-----------
\section{教育背景}
\resumeSubHeadingListStart
\resumeSubheading
{\href{https://hku.hk/}{香港大学}}{香港}
\resumeSubSubheading
{计算机科学理学硕士（研究生）}{2025.09 -- 2026.11 \scriptsize(预计)}
\resumeItemListStart
\resumeItem{GPA: 3.39/4.3；相关课程：量子计算、数据挖掘、算法交易、信息安全、大数据管理、UI/UX 开发}
\resumeItemListEnd
\vspace{-3pt}
\resumeSubSubheading
{计算机科学工学学士（本科）}{2021.09 -- 2025.06}
\resumeItemListStart
\resumeItem{GPA: 3.56/4.3，一等荣誉学位，院长荣誉名单；相关课程：数据结构、算法、软件工程、数据库、网络通信、操作系统、计算机组成原理、机器学习、计算机视觉、自然语言处理、深度学习、数据科学、面向对象程序设计}
% \resumeItem{相关课程：机器学习、计算机视觉、操作系统、面向对象程序设计、数据科学、软件工程、算法与数据结构、数据库、网络通信、自然语言处理、深度学习}
\resumeItemListEnd
\resumeSubHeadingListEnd

% MARK: EXPERIENCE
%-----------EXPERIENCE-----------
\section{实习经历}
\resumeSubHeadingListStart

\resumeSubheading
{\href{https://www.apple.com/}{苹果公司}}{上海}
\resumeSubSubheading
{AI 开发实习生}{2026.05 -- 2026.09}
\resumeItemListStart
\resumeItem{主导测试工程师故障监控分析 Agent 插件的设计与交付：完成痛点访谈、架构设计，封装为 Agent SKILL，配套开发 MCP Server、CLI 工具、后台系统，并整合 Human-in-the-Loop 工单系统，支持人工一键创建与关联工单。}
\resumeItem{针对海量测试数据场景，设计定制 MCP Server 将数据拉取至本地，由脚本完成计算与汇总，全程不将原始数据载入 Agent 上下文，规避大模型在大规模数值计算上的短板；按站点拆分任务，由多个子 Agent 并行处理各自站点的错误信息分析与归因，实现每日自动化报告生成，显著缩短工程师制作报告的时间，助力问题提前发现与解决。}
\resumeItem{搭建并完善 Skill 仓库的 CI/CD 发布流程，包括 AI PR 检查、自动打包制品发布、一行命令安装及智能体自更新机制。}
\resumeItemListEnd

\vspace{4pt}

\resumeSubheading
{\href{https://www.hku.hk/}{香港大学}}{香港}
\resumeSubSubheading
{研究助理 (兼职)}{2025.07 -- 至今}
\resumeItemListStart
\resumeItem{针对港大法律与科技中心的需求开发适配教学的法律AI工具，负责工具的研发与落地，包括对香港原始法律文书做结构化改造，并搭建关系型语义数据库，为Agentic RAG大语言模型法律应用提供高效信息检索支撑。}
\resumeItem{负责维护并持续优化 \href{https://hklii.hk/}{\thinunderline{HKLII}} 平台的 AI \href{https://ai.hklii.hk/dt-predictor/}{\thinunderline{量刑预测器}}，
利用大模型对判决书进行结构化信息提取，搭建人工检查平台，并对数据进行分析及预测模型搭建，使系统及时适配香港最新法律条文。\footnotesize (\href{https://github.com/Xiang-CH/hklii-drug-trafficking-sentence-predictor}{\thinunderline{项目仓库}})}
\resumeItemListEnd

\resumeSubSubheading
{学生研究助理}{2023.06 -- 2024.05，2024.09 -- 2025.07}
\resumeItemListStart
\resumeItem{在香港大学 Innovation Wing 设计并主导多场实践型 AI 工作坊，内容包括基于 Azure AI 服务构建自定义聊天机器人，以及结合版面检测（Layout Detection）的多模态 RAG 文档理解应用开发等。}
\resumeItem{参与并主导多个跨部门学生 AI 项目，包括 \href{https://clic-search.vercel.app/}{\thinunderline{CLIC-Search}}（社区法治信息中心语义搜索引擎）与多模态历史研究助手。}
\resumeItemListEnd

\vspace{4pt}

\resumeSubheading
{\href{https://www.iflytek.com/}{科大讯飞}}{深圳}
\resumeSubSubheading
{软件开发助理工程师}{2024.06 -- 2024.08}
\resumeItemListStart
\resumeItem{参与教育产品线 AI 英语口语私教与 AI 绘本陪读 等 LLM 应用的敏捷开发，使用 \emph{Node-RED} 设计并维护应用工作流；并通过 LoRA 微调优化模型能力，以满足产品的定制化功能需求。}
\resumeItem{通过 \emph{Python} 脚本自动化模型微调所需的数据工程流程，并对提示词调优流程进行优化以提升效率。}
\resumeItemListEnd

\resumeSubHeadingListEnd

% MARK: PROJECTS
%-----------PROJECTS-----------
\section{项目经历}
\resumeSubHeadingListStart
\resumeProjectHeading
{\textbf{\href{https://clic.cxiang.site/search}{\thinunderline{CLIC-Chat}}} $|$ \emph{Next.js, Prisma, MS-SQL, Azure AI, LLM, RAG}}{2025.07 -- 至今}
\resumeItemListStart
\resumeItem{开发面向香港法律内容问答的 全栈 AI Web 应用，使用 \emph{Next.js} 与 \emph{Prisma} 构建系统架构。}
\resumeItem{解析并清洗香港法律数据源（包括法规、判决书及法律文章），提取结构化信息并构建基于 \emph{MS-SQL} 的关系型数据库。}
\resumeItem{通过对非结构化法律文本进行语义分块与向量化，在 \emph{Azure AI Search} 上构建 向量 + 关键词（BM25）混合检索索引。}
\resumeItem{实现基于 RAG 的问答系统，包括查询扩展（Query Expansion）、混合检索、关系信息提取与结果重排序（Reranking），以提升检索相关性并生成更准确的 LLM 回答。}
\resumeItemListEnd

% \resumeProjectHeading
% {\textbf{\href{https://cucampus.one/}{\thinunderline{CU-Campus}}} $|$ \emph{React Router, Tailwind CSS, Vite, i18n}}{2024.05 -- 2024.10}
% \resumeItemListStart
% \resumeItem{CU Campus 是面向香港中文大学学生的课程评价与课表排期平台，整合多项校园服务并提供一站式信息入口。}
% \resumeItem{参与产品 UX 设计、功能规划及敏捷迭代开发。}
% \resumeItem{基于 \emph{React Router} 实现前端开发，使用 \emph{i18n} 完成三语本地化，并与后端 RESTful API 进行数据交互。}
% \resumeItemListEnd

\resumeProjectHeading
{更多项目：\href{https://cxiang.site/project}{\textbf{\thinunderline{cxiang.site/project}}}}{}
\resumeSubHeadingListEnd



% MARK: SKILLS
%-----------PROGRAMMING SKILLS-----------
\section{技术栈}
\begin{itemize}[leftmargin=0.15in, label={}]
	\fontsize{10pt}{11.5pt}\selectfont\color[rgb]{0.1,0.1,0.1}{\item{
	      \textbf{编程语言}{: TypeScript, JavaScript, Python, SQL, HTML/CSS} \\
	      \textbf{框架与平台}{: Node.js, Next.js, Tanstack, Flask, FastAPI, Tailwind CSS, Prisma, Chrome 插件, 微信小程序} \\
	      \textbf{数据库}{: MS-SQL, MongoDB, PostgreSQL, MySQL, ChromaDB} \\
	      \textbf{开发工具}{: Git, Docker, Linux, Azure, AWS, 阿里云, Postman, Cloudflare, Vercel, Dify, n8n, Codex} \\
	      \textbf{常用库}{: React, Pandas, NumPy, Matplotlib, PyTorch, Scikit-learn, LangChain, Selenium, Vercel AI SDK}
	      }}
\end{itemize}


%-------------------------------------------
\end{document}
